\documentclass{jsarticle}
\usepackage[dvipdfmx,final]{graphicx}
\usepackage{amsmath}
\usepackage{amssymb}
\begin{document}
\title{Gamma function relativity restream into Mebius \\
Klein dimension, Cohomological define with Euler product, \\
Euler number from integral manifold to general relativity \\
and Special relativity. Imaginary pole of Volume manifold \\
emerged with Beta function.}
\author{Masaaki Yamaguchi}
\date{}
\maketitle

        オイラー数を多様体積分で施すと、
   ガンマ関数になる。この場合は、Frobeniusの定理によって、
   和と積の法則、確率論から導かれる結果、
   異次元と宇宙の関係に持っていける。
   オイラー数を多様体積分すると、ガンマ関数になる。
   オイラー数における頂点が、宇宙の特異点に対応して、
   辺がノルムに属して、面が一般相対性理論における無限の接線であり、
   $$ \text{F(頂点) - N(辺) + E(面)} = 2 $$
   $$ {(x,y,z) \over \Gamma} = \int{(- F\circ N + E)} + 
   (-F\circ E + N) + (N\circ E + F)dx_m $$
   $$ {{{\partial \over \partial f_S}(x,y,z)} \over {\Gamma}} $$
   ここで、オイラーの定数をも多様体積分にして、付け加えると、
   $$ = \int{({\int {1 \over x^s}dx} - \log x)}\mathrm{dvol} $$
   大域的多様体積分のガンマ関数に化ける。
   $$ = \int e^{-x}x^{1 - t}dx_m $$

   $AdS_5$多様体の方程式を、オイラー数の各部分単体に合わせると、
   $$ ||ds^2|| = e^{-2\pi T||\psi||}[\eta + \bar{h}(x)]dx^{\mu}dx^{\nu} 
   + T^2d^2 \psi $$
   開集合として、表せる。
   $$ = \mathcal{O}(x) $$
   この上の式たちから、ダランベルシアンは、$\square$と、オイラー数から演算子
   が決まっていて、共変微分は、$\nabla$で、双対被覆での単体での演算子で、
   決まっている。
   $$ \square = {{8 \pi G} \over c^4} T^{\mu\nu} $$
   $$ \nabla \psi = 4 \pi G \rho $$

   オイラー数を共変微分すると、偏微分方程式と同じ機能を、
   大域的積分多様体のオイラー数が、受け持っていて、
   このオイラー数の大域的多様体の共変微分で、ガンマ関数の大域的多様体
   が、
   ここで、オイラーの定数をも多様体積分にして、付け加えると、
   $$ = \int{({\int {1 \over x^s}dx} - \log x)}\mathrm{dvol} $$
   大域的多様体積分のガンマ関数に化ける。
   $$ = \int e^{-x}x^{1 - t}dx_m $$
   $$ = {d \over d f}F = \int \Gamma(\gamma)^{'}dx_m $$
   $$ = \int C dx_m = \int \Gamma dx_m \circ {d \over d \gamma}\Gamma $$
   $$ \le \int \Gamma dx_m + {d \over d \gamma}\Gamma $$
   $$ = e^{- \theta} + e^{i \theta} $$
   $$ = \square = 2(\cos (i x \log x) - i \sin (i x \log x))$$
   $$ = e^{-f} - e^{f} \le e^{-f} + e^{f} $$
   統一場理論は、彩さんのJones多項式に行き着く。



\end{document}
